% Chapter 2 -- Getting Started.
% Sources (Law 1): os/milton/command.c (builtin_dir/type/copy/del/cd/md/rd,
% the REPL prompt $P$G, read_line CTRL-C discard, AUTOEXEC.BAT probe);
% os/milton/kmain.c (start-up order, FAT-MOUNT root listing, SYSINIT
% CONFIG.SYS-or-standard settings); spec/dos_messages.json;
% spec/dos_banner.txt; pubs/screens/disk/MANIFEST (the practice diskette).
% Every figure is a capture of the shipped COMMAND.COM image
% (build/tracer_boot.img) driven by screens/SCREENS through
% pubs/screens/capture.sh; `make test-pubs-screens` re-captures them.
% GATES: test-boot test-shell test-dir test-command test-copy-selfcopy test-dos-safety-k01 test-dos-safety-k03 test-dos-safety-k04 test-j1lu-control-emu test-ut6d test-fat-subdir test-pubs-screens
\Gates{test-boot, test-shell, test-dir, test-command, test-copy-selfcopy,
  test-dos-safety-k01, test-dos-safety-k03, test-dos-safety-k04,
  test-j1lu-control-emu, test-ut6d, test-fat-subdir, test-pubs-screens}
\chapter{Getting Started}

This chapter is a series of short exercises.  Do them in order, at the
computer, with this book open beside you.  Each exercise shows the screen
as it should look when you have finished.

\section{Before You Begin}
The exercises use the InitechDOS practice diskette.  It holds three text
files, \file{README.TXT}, \file{MEMO.TXT} and \file{PHONE.TXT}; the file
\file{AUTOEXEC.BAT}; the directory \file{APPS}, which holds the utility
programs; and the directory \file{SAMPLES}, which holds the batch files used
in Chapter~5.  You will copy, delete and rename files, so use a copy of the
practice diskette, not the original.
\draft{bead needed: the package and the diskette-copying command (DISKCOPY)
are not built; the screens use a practice volume minted by
pubs/screens/capture.sh as drive A}

\section{Starting the System}
\index{starting InitechDOS}
\begin{Steps}
\item Make sure the practice diskette is in place as drive~A.
\item Turn on the display, then the computer.
\end{Steps}
The computer checks its memory and then loads InitechDOS.  This takes a few
seconds.  When InitechDOS is ready, the screen looks like
Figure~\ref{fig:startup}.

\Screen[rows=8, label=fig:startup]{screens/startup.txt}{The start-up screen}

\SideHead{Copyright notice}The first two lines show the version of
InitechDOS and the copyright notice.

\SideHead{Start-up}\draft{initech-qskz: the shipped build lists the files in
the root directory of drive A here (Figure~\ref{fig:startup}); DOS 3.3 asks for the date
and time when there is no AUTOEXEC.BAT}When the disk has no
\file{AUTOEXEC.BAT} file, InitechDOS asks for the date and the time; press
\key{Enter} at each question to keep the date and time shown.

\SideHead{Command prompt}The last line, \file{A:\bs>}, is the command
prompt.  It tells you that InitechDOS is waiting for a command.

\SideHead{Configuration}Before the prompt appears, InitechDOS reads
\file{CONFIG.SYS} from drive~A to learn how your system is set up; if there
is no such file, it uses the standard Initech settings (Chapter~6).  Then
the command processor carries out the commands in \file{AUTOEXEC.BAT}.  The
practice diskette's \file{AUTOEXEC.BAT} tells InitechDOS where to find the
utility programs, and displays nothing (Chapter~5).

\section{The Command Prompt}
\index{command prompt}\index{prompt}
The command prompt shows the current drive and directory, followed by the
character \file{>}.  At \file{A:\bs>} the current drive is drive~A and the
current directory is its root directory.  When you change to another
directory, the prompt changes to match, as you will see later in this
chapter.  You can change the form of the prompt with the \cmd{PROMPT}
command (Chapter~6); this book always shows the standard prompt.

\section{Typing Commands}
\index{typing commands}
Type a command at the prompt, then press \key{Enter}.  InitechDOS does
nothing with the command until you press \key{Enter}, so you can correct
mistakes first:
\begin{itemize}
\item Press \key{Backspace} to erase the character to the left of the
  cursor.\index{Backspace key}
\item Press \key{Ctrl-C} to cancel the whole line.  InitechDOS ignores what
  you typed and displays the prompt again.\index{Ctrl-C}
\end{itemize}
Chapter~8 describes these keys.  You can type commands in capital or small
letters.  This book shows them in capitals.  If InitechDOS does not
recognize a command, or a command is missing information it needs, it
displays a message and the prompt again.  Nothing is changed on the disk.
Try it now:
\begin{Steps}
\item Type \type{CLS} and press \key{Enter}.  The \cmd{CLS}%
  \index{CLS@\texttt{CLS}} command clears the screen and displays the prompt
  at the top.
\item Type \type{B:} and press \key{Enter}.
\item Type \type{DIRECTORY} and press \key{Enter}.
\item Type \type{TYPE} and press \key{Enter}.
\end{Steps}
Your screen now looks like Figure~\ref{fig:errors}.

\Screen[rows=7, label=fig:errors]{screens/errors.txt}{Two kinds of mistake}

\cmd{B:} and \cmd{DIRECTORY} are not commands that InitechDOS knows, so it
answers \texttt{Bad command or file name}.  \cmd{TYPE} is a command, but it
needs the name of a file to work on; \texttt{Required parameter missing}
tells you that something is missing from the command line.

\section{Listing the Files on a Disk}
\index{listing files}\index{DIR@\texttt{DIR}}
The \cmd{DIR} command lists the files in the current directory.
\begin{Steps}
\item Type \type{CLS} and press \key{Enter} to clear the screen.
\item Type \type{DIR} and press \key{Enter}.
\end{Steps}

\Screen[rows=15, label=fig:dir]{screens/dir.txt}{A directory listing}

The listing in Figure~\ref{fig:dir} begins with the volume label of the disk
and the name of the directory being listed.  Then comes one line for each
file or directory:

\SideHead{Name}The name of the file, followed by its extension in a column of
its own.  \file{README TXT} is the file \file{README.TXT}.

\SideHead{Size}The size of the file in bytes.  A directory shows
\texttt{<DIR>} instead.

\SideHead{Date and time}When the file was last changed: month, day and
year, then the time of day, with \file{a} for morning and \file{p} for
afternoon.

The last line gives the number of entries listed and the number of bytes
still free on the disk.

\section{Displaying a File}
\index{displaying a file}\index{TYPE@\texttt{TYPE}}
The \cmd{TYPE} command displays the contents of a text file on the
screen.
\begin{Steps}
\item Type \type{CLS} and press \key{Enter}.
\item Type \type{TYPE MEMO.TXT} and press \key{Enter}.
\end{Steps}

\Screen[rows=10, label=fig:type]{screens/type.txt}{Displaying MEMO.TXT}

\cmd{TYPE} shows the file exactly as it is stored.  Use it for text files
only: a program file contains codes that are not meant to be displayed and
will fill the screen with strange characters.

\section{Copying a File}
\index{copying files}\index{COPY@\texttt{COPY}}
Make a copy of a file before you change it, so that you can go back to the
original.  The \cmd{COPY} command copies one file to another.  Name the file
to copy from first, and the file to copy to second.
\begin{Steps}
\item Type \type{CLS} and press \key{Enter}.
\item Type \type{COPY MEMO.TXT MEMO.BAK} and press \key{Enter}.
\item Type \type{DIR} and press \key{Enter} to see the new file.
\end{Steps}

\Screen[rows=18, label=fig:copy]{screens/copy.txt}{Copying MEMO.TXT to MEMO.BAK}

\cmd{COPY} reports \texttt{1 file(s) copied}, and the directory now lists
\file{MEMO.BAK}, the same size as \file{MEMO.TXT}.
\draft{initech-dwbl: the shipped build stamps the new file 0-00-80 12:00a
(Figure~\ref{fig:copy}); DOS 3.3 gives a copy the date and time of the original}
The copy has the same date and time as the original file.

If the file you copy to already exists, \cmd{COPY} replaces it.  If you
name the same file twice, \cmd{COPY} refuses with \texttt{File cannot be
copied onto itself} and copies nothing.  \cmd{COPY} copies one file at a
time; to copy a group of files, use \cmd{XCOPY} (Chapter~9).

\section{Deleting Files}
\index{deleting files}\index{DEL@\texttt{DEL}}
The \cmd{DEL} command removes a file from the disk.  \cmd{ERASE} is another
name for the same command.  For this exercise a copy, \file{MEMO.BAK}, has
been made first.
\begin{Steps}
\item Type \type{CLS} and press \key{Enter}.
\item Type \type{DEL MEMO.BAK} and press \key{Enter}.
\item Type \type{TYPE MEMO.BAK} and press \key{Enter}.
\end{Steps}

\Screen[rows=4, label=fig:del]{screens/del.txt}{Deleting MEMO.BAK}

\cmd{DEL} displays nothing when it succeeds.  The \cmd{TYPE} command shows
that the file is gone.  \cmd{DEL} accepts the wildcard names described in
Chapter~3, and asks \texttt{Are you sure (Y/N)?} before it deletes every
file in a directory.

\begin{Important}
A deleted file cannot be brought back with InitechDOS commands.  Check the
name before you press \key{Enter}.
\end{Important}

\section{Working with Directories}
\index{directories}
In this exercise you create a subdirectory, copy a file into it, and look at
its contents.  Three commands are used:

\SideHead{MD}\cmd{MD} (make directory) creates a subdirectory.
\cmd{MKDIR} is another name for it.\index{MD@\texttt{MD}}

\SideHead{CD}\cmd{CD} (change directory) makes another directory the current
one.  \cmd{CHDIR} is another name for it.  Typed alone, \cmd{CD} displays
the current directory.\index{CD@\texttt{CD}}

\SideHead{RD}\cmd{RD} (remove directory) removes an empty subdirectory.
\cmd{RMDIR} is another name for it.\index{RD@\texttt{RD}}

\begin{Steps}
\item Type \type{CLS} and press \key{Enter}.
\item Type \type{MD REPORTS} and press \key{Enter}.
\item Type \type{COPY MEMO.TXT REPORTS} and press \key{Enter}.
\item Type \type{CD REPORTS} and press \key{Enter}.  The prompt changes to
  \file{A:\bs REPORTS>}.
\item Type \type{DIR} and press \key{Enter}.
\end{Steps}

\Screen[rows=16, label=fig:paths]{screens/paths.txt}{The REPORTS subdirectory}

Every subdirectory contains two special entries.  \file{.} stands for the
directory itself, and \file{..} for the directory that contains it, its
\term{parent}.  To go back up to the parent directory, type \type{CD ..}.
\draft{initech-dwbl: the new directory and its copy are stamped 0-00-80
12:00a in the shipped build (Figure~\ref{fig:paths}); DOS 3.3 stamps the current date
and time}

To finish the exercise, remove the subdirectory.  A directory must be empty
before \cmd{RD} will remove it, so delete its file first:
\begin{TypeIn}
CD ..
DEL REPORTS\MEMO.TXT
RD REPORTS
\end{TypeIn}
If the directory still contains files, \cmd{RD} displays \texttt{Invalid
path, not directory, or directory not empty} and leaves it in place.

\section{Where to Go Next}
You can now start InitechDOS, list and display files, copy and delete
them, and organize them in directories.  Chapter~3 tells you more about
file names and directories; the \doctitle{InitechDOS Reference} describes
every command in detail.
